\section{Conclusion}

Accessibility in Jetpack Compose is not a checkbox but a design discipline that spans the whole
stack. This project took an ordinary podcast UI and made it usable without sight by working the
semantics tree deliberately at every level: naming and hiding elements correctly
(Section~\ref{sec:core}), then shaping how those elements are grouped, ordered, and acted upon
(Section~\ref{sec:adv}). The result is a three-screen app in which a TalkBack user can discover a
show, choose an episode, and control playback --- hearing a clean sentence per card, reaching every
action through the actions menu, and getting spoken feedback the moment playback state changes.

Equally important is the evaluation discipline of Section~\ref{sec:test}: TalkBack answers the
listening question, Scanner triages quickly, Layout Inspector explains the structure, Compose UI
Testing protects semantics against regressions, and ATF automates the mechanical checks. The value
lies in \emph{combining} these layers, defining a clear scope, and --- as our Agoda case study
showed --- never mistaking a tool's raw suggestion for a confirmed defect.

Our broader takeaway is that accessible design is cheaper and better when it is built in from the
first composable rather than retrofitted: preferring Material components with sound default
semantics, reaching for custom semantics only where the design demands it, and testing with the same
assistive technology real users rely on. Building inclusively is, in the end, simply building well.

\subsection*{Future work}

The natural Version~2 is to run TalkBack on our own build across a device/version matrix, capture
Layout Inspector evidence for representative findings, expand the Compose test suite and wire ATF
into CI, and --- resources permitting --- invite users who rely on assistive technology to evaluate
the core tasks.

\section*{Appendix on AI usage}
\addcontentsline{toc}{section}{Appendix on AI usage}

AI assistants were used as a drafting and formatting aid: to help structure this report, generate the
LaTeX scaffolding and diagrams, and produce the annotated UI mockups from our own designs and code.
All accessibility techniques, the app architecture, the research decisions, and the case-study data
originate from the group's own work; every AI-assisted passage was reviewed and edited by the members
responsible for that section. Code samples in this report are excerpts from the app we wrote, not
AI-generated boilerplate.
